\section{Experiments}
\label{sec:experiments}

\subsection{Experimental Setup}
\label{sec:setup}

\paragraph{Datasets.} We evaluate on two bearing run-to-failure datasets with complementary characteristics:

\begin{table}[t]
\centering
\caption{Dataset summary. Episodes are complete run-to-failure sequences from initial operation to bearing failure.}
\label{tab:datasets}
\small
\begin{tabular}{@{}lcccccc@{}}
\toprule
Dataset & Episodes & Sampling Rate & Channels & Lifetime Range & Operating Conditions \\
\midrule
FEMTO (PRONOSTIA) & 16 & 25.6\,kHz & 1 (vibration) & 1--7\,h & 3 load/speed settings \\
XJTU-SY & 7 & 25.6\,kHz & 2 (H+V vib.) & 0.5--2.6\,h & 3 load/speed settings \\
\midrule
\textbf{Combined} & \textbf{23} & --- & --- & \textbf{CV=0.635} & \textbf{6 conditions} \\
\bottomrule
\end{tabular}
\end{table}

\plannedc{
\begin{table}[t]
\centering
\caption{C-MAPSS turbofan engine dataset summary.}
\label{tab:cmapss}
\small
\begin{tabular}{@{}lccccc@{}}
\toprule
Subset & Train engines & Test engines & Operating conditions & Fault modes & Sensors \\
\midrule
FD001 & 100 & 100 & 1 & 1 (HPC) & 21 \\
FD002 & 260 & 259 & 6 & 1 (HPC) & 21 \\
FD003 & 100 & 100 & 1 & 2 (HPC+Fan) & 21 \\
FD004 & 249 & 248 & 6 & 2 (HPC+Fan) & 21 \\
\bottomrule
\end{tabular}
\end{table}
}

\paragraph{Evaluation protocol.} All methods are evaluated with 75\%/25\% episode-based train/test splits, fixed across methods. We report RMSE on normalized RUL ($\in [0, 1]$) averaged over 5 seeds (10 seeds for cross-dataset experiments) with standard deviations. Statistical significance is assessed via paired $t$-tests.

\paragraph{Baselines.} We compare 11 methods spanning naive baselines, end-to-end deep learning, handcrafted features, and self-supervised approaches:
\begin{itemize}[nosep]
  \item \textbf{Naive}: Constant mean predictor, elapsed-time-only baseline.
  \item \textbf{Signal processing}: Envelope RMS + LSTM.
  \item \textbf{End-to-end}: CNN-LSTM, CNN-GRU-MHA.
  \item \textbf{Handcrafted}: HC+LSTM, HC+MLP, Transformer+HC.
  \item \textbf{Self-supervised}: Random JEPA+LSTM (untrained encoder), JEPA+LSTM (pretrained).
  \item \textbf{Hybrid}: JEPA+HC Transformer (our best).
\end{itemize}

\subsection{In-Domain RUL Prediction}
\label{sec:indomain}

\begin{table}[t]
\centering
\caption{In-domain RUL prediction results (mixed FEMTO+XJTU-SY, 5-seed average). Improvement is relative reduction in RMSE vs.\ elapsed-time-only baseline (RMSE=0.224).}
\label{tab:main_results}
\small
\begin{tabular}{@{}llccc@{}}
\toprule
Category & Method & RMSE & $\pm$std & vs.\ Time-Only \\
\midrule
\multirow{2}{*}{Naive} & Constant mean & 0.290 & --- & $-29.4\%$ \\
& Elapsed time only & 0.224 & --- & 0\% \\
\midrule
Signal proc. & Envelope RMS + LSTM & 0.287 & 0.001 & $-28.1\%$ \\
\midrule
\multirow{2}{*}{End-to-end} & CNN-LSTM & 0.195 & 0.005 & $+13.0\%$ \\
& CNN-GRU-MHA & 0.185 & 0.005 & $+17.4\%$ \\
\midrule
\multirow{3}{*}{Handcrafted} & HC + LSTM & 0.177 & 0.016 & $+21.2\%$ \\
& HC + MLP & 0.085 & 0.004 & $+62.0\%$ \\
& Transformer + HC & 0.070 & 0.006 & $+68.9\%$ \\
\midrule
\multirow{2}{*}{Self-supervised} & Random JEPA + LSTM & 0.221 & 0.008 & $+1.5\%$ \\
& JEPA + LSTM & 0.189 & 0.015 & $+15.8\%$ \\
\midrule
\textbf{Hybrid} & \textbf{JEPA + HC Transformer} & \textbf{0.055} & \textbf{0.004} & $\mathbf{+75.5\%}$ \\
\bottomrule
\end{tabular}
\end{table}

\Cref{tab:main_results} presents the complete comparison. Several findings emerge:

\paragraph{JEPA pretraining is effective.} JEPA+LSTM (RMSE 0.189) significantly outperforms both the elapsed-time baseline ($+15.8\%$, $p = 0.010$) and random encoder initialization ($p = 0.010$). The pretrained encoder learns meaningful vibration representations without any RUL labels.

\paragraph{JEPA matches handcrafted features.} The difference between JEPA+LSTM (0.189) and HC+LSTM (0.177) is not statistically significant ($p = 0.40$), demonstrating that self-supervised pretraining can match decades of domain expertise in feature engineering---without requiring that expertise.

\paragraph{The hybrid wins decisively.} Combining JEPA embeddings with handcrafted features in a Transformer architecture achieves the best result: RMSE $0.055 \pm 0.004$, a $75.5\%$ improvement over elapsed-time and $20.7\%$ over the best handcrafted-only method (Transformer+HC at 0.070). This confirms that JEPA captures complementary information to spectral statistics.

\subsection{Cross-Dataset Transfer}
\label{sec:transfer}

Cross-dataset transfer---training on one bearing test rig and evaluating on another---is the critical test of learned representations. Different test rigs have different operating conditions, bearing types, and lifetime distributions (coefficient of variation = 0.635).

\begin{table}[t]
\centering
\caption{Cross-dataset transfer results (10-seed average). Elapsed time baseline degrades sharply cross-dataset due to lifetime distribution mismatch.}
\label{tab:transfer}
\small
\begin{tabular}{@{}lcccc@{}}
\toprule
Configuration & Elapsed Time & JEPA+LSTM & Contrastive+LSTM & Contrastive vs.\ JEPA \\
\midrule
FEMTO$\to$FEMTO (within) & 0.027 & $0.113 \pm 0.011$ & $0.142 \pm 0.012$ & $-20.4\%$ \\
XJTU$\to$XJTU (within) & 0.159 & 0.195 & 0.214 & $-9.7\%$ \\
\midrule
\textbf{FEMTO$\to$XJTU (cross)} & \textbf{0.367} & $0.280 \pm 0.007$ & $\mathbf{0.227 \pm 0.015}$ & $+18.8\%$ ($p < 0.001$) \\
\textbf{XJTU$\to$FEMTO (cross)} & \textbf{0.336} & 0.403 & $\mathbf{0.309 \pm 0.007}$ & $+23.2\%$ \\
\bottomrule
\end{tabular}
\end{table}

\paragraph{Contrastive dominates cross-dataset.} Temporal contrastive learning significantly outperforms JEPA for cross-dataset transfer ($+18.8\%$ on FEMTO$\to$XJTU, $p < 0.001$). This is because the contrastive objective forces the encoder to learn what \emph{changes} over a bearing's life---the spectral centroid shift---which is the invariant degradation signature across datasets.

\paragraph{JEPA wins within-dataset.} Conversely, JEPA achieves lower RMSE within datasets, especially on FEMTO (0.113 vs 0.142). JEPA's patch prediction objective preserves fine-grained waveform details that are useful when train and test distributions match.

\paragraph{Both beat elapsed time.} In the cross-dataset setting, where elapsed time baselines degrade sharply (RMSE 0.336--0.367), both learned representations provide substantial improvement ($+24\%$ for JEPA, $+38\%$ for contrastive vs time-only).

\plannedc{
\subsection{Multivariate Input Ablation}
\label{sec:multivariate_exp}

\begin{table}[t]
\centering
\caption{Effect of input channels on RUL prediction (Paderborn dataset).}
\label{tab:multivariate}
\small
\begin{tabular}{@{}lccc@{}}
\toprule
Input Configuration & Channels & RMSE ($\pm$std) & vs.\ 1-channel \\
\midrule
Single vibration & 1 & \placeholder{---} & --- \\
Dual vibration (radial + axial) & 2 & \placeholder{---} & \placeholder{---} \\
All vibration channels & 4 & \placeholder{---} & \placeholder{---} \\
Full multivariate & 8 & \placeholder{---} & \placeholder{---} \\
\bottomrule
\end{tabular}
\end{table}

\placeholder{Discussion of multivariate results, which channels contribute most, relationship between physics groups and predictive value.}
}

\plannedc{
\subsection{Spatiotemporal Masking Ablation}
\label{sec:masking_exp}

\begin{table}[t]
\centering
\caption{Effect of masking strategy during JEPA pretraining on downstream RUL prediction.}
\label{tab:masking}
\small
\begin{tabular}{@{}lccc@{}}
\toprule
Masking Strategy & Description & RMSE ($\pm$std) & vs.\ Random \\
\midrule
Random & Independent per patch/channel & \placeholder{---} & --- \\
Physics-group & Mask all channels in a physics group & \placeholder{---} & \placeholder{---} \\
Spatiotemporal & Block masking in time $\times$ channel & \placeholder{---} & \placeholder{---} \\
Curriculum & Random $\to$ spatiotemporal & \placeholder{---} & \placeholder{---} \\
\bottomrule
\end{tabular}
\end{table}

\placeholder{Discussion of masking strategy results, whether physics-aware masking forces better cross-modal representations.}
}

\plannedc{
\subsection{C-MAPSS Turbofan Benchmark}
\label{sec:cmapss}

To validate generalization beyond bearings, we evaluate on the C-MAPSS turbofan engine degradation dataset~\citep{saxena2008cmapss}.

\begin{table}[t]
\centering
\caption{C-MAPSS results (Score metric, lower is better). Published baselines from respective papers.}
\label{tab:cmapss_results}
\small
\begin{tabular}{@{}lcccc@{}}
\toprule
Method & FD001 & FD002 & FD003 & FD004 \\
\midrule
DCNN~\citep{li2018dcnn} & 274 & 10,412 & 284 & 12,466 \\
LSTM~\citep{zheng2017lstm} & 338 & 4,450 & 852 & 5,550 \\
TCN-Transformer & \placeholder{---} & \placeholder{---} & \placeholder{---} & \placeholder{---} \\
\midrule
Ours (JEPA+LSTM) & \placeholder{---} & \placeholder{---} & \placeholder{---} & \placeholder{---} \\
Ours (Contrastive+LSTM) & \placeholder{---} & \placeholder{---} & \placeholder{---} & \placeholder{---} \\
Ours (Hybrid) & \placeholder{---} & \placeholder{---} & \placeholder{---} & \placeholder{---} \\
\bottomrule
\end{tabular}
\end{table}

\placeholder{Discussion of C-MAPSS results, comparison with published SOTA, analysis of which subsets benefit most from self-supervised pretraining.}
}

\plannedc{
\subsection{Synthetic Grey Swan Augmentation}
\label{sec:augmentation_exp}

\begin{table}[t]
\centering
\caption{Effect of synthetic augmentation on prediction of rare failure modes.}
\label{tab:augmentation}
\small
\begin{tabular}{@{}lcc@{}}
\toprule
Training Data & RMSE (common modes) & RMSE (rare modes) \\
\midrule
Real failures only & \placeholder{---} & \placeholder{---} \\
Real + synthetic augmentation & \placeholder{---} & \placeholder{---} \\
\bottomrule
\end{tabular}
\end{table}

\placeholder{Discussion of augmentation results, which rare modes benefit most, qualitative comparison of synthetic vs real degradation signatures.}
}
